\section{Conclusion}\label{sec:conclusion}

\subsection{Stable closures develop something price-like}

This paper began from a familiar but under-explained recurrence: mature theories repeatedly introduce small sets of variables that budget feasibility, accumulate along protocols, or appear as exchange rates between incompatible aims. We argued that these variables should be treated neither as loose metaphors nor as merely local optimization artifacts. Within the Six Birds closure calculus, they are currencies: low-dimensional structural coordinates that govern what a layer can afford in order to stay closed.

The paper's main principle is therefore cross-layer. A lower-layer currency becomes a higher-layer constraint because packaged higher-level objects can persist only under bounded lower-layer spending. Once the higher layer is organized by those inherited budgets, its own currency emerges as the corresponding shadow price. The recurrence of temperature, price, regularization, and attention-like tradeoff parameters is thus not accidental. It is the signature of closure living on a finite substrate.

The finite-state realization made this principle measurable. Five signatures supported the account: honest path-audit monotonicity under coarse observation, monotone price emergence with a slack regime, growth of currency dimension with resolution, predictive and inferential failure of proxy currencies, and reduced packaging defect under stronger budget enforcement. A Lean formalization supplied a narrower but important formal anchor by proving deterministic-pushforward monotonicity for a finite KL quantity.

The broader claim is modest and strong at the same time. It is modest because the notion of currency is layer-relative, the laboratory is finite-state, and the formal proof anchors only one part of the story. It is strong because the same mathematical pattern appears across physics, information, economics, learning, and cognition whenever stable description must be maintained under bounded spending. Stable closures do not merely accumulate description. They develop budgets, and budgets develop prices. Stable closures develop something price-like.
